\begin{textsample}{POS Dim 3 – mt – Score 16.00 – mt/mt\_inf\_14.txt}  \label{ex:f3_pos_090}
2.1.4 As possibilidades de experiências com \textbf{Bebês}, \textbf{Crianças} Bem Pequenas e \textbf{Crianças} Pequenas² Na Educação \textbf{Infantil}, o trabalho pedagógico pautado nos campos de experiências precisa contemplar os eixos norteadores interações e \textbf{brincadeiras}, as dez competências gerais e os direitos e os objetivos de aprendizagem e desenvolvimento, para cada grupo etário mencionado na BNCC ( 2017 ).
Nessa perspectiva e considerando as especificidades do Estado de Mato Grosso, faz -se necessário, dentre as diversas possibilidades de ações pedagógicas com \textbf{bebês}, \textbf{crianças} bem \textbf{pequenas} e \textbf{crianças} \textbf{pequenas}, propiciar experiências de : - Criar situações em que as \textbf{crianças} expressem seus \textbf{afetos}, \textbf{desejos}, saberes e aprendam a ouvir o outro, a \textbf{conversar} e negociar argumentos e metas, a fazer planos comuns, a enfrentar conflitos, a participar de uma atividade do grupo, a criar amizades com seus companheiros ; - \textbf{Apoiar} as \textbf{crianças} a desenvolver uma identidade pessoal, um sentimento de autoestima, autonomia, confiança em suas possibilidades e de pertencimento a um determinado grupo étnico-racial, crença religiosa, local de nascimento, etc. ; - Fortalecer os vínculos afetivos de todas as \textbf{crianças} com suas famílias e ajudá -las a captar as possibilidades trazidas por diferentes tradições culturais para a compreensão do mundo e de si mesmas ; - Construir com as \textbf{crianças} o entendimento da importância de cuidar de sua saúde e bem-estar, no decorrer das atividades cotidianas ; - Criar com as \textbf{crianças} \textbf{hábitos} ligados à \textbf{limpeza} e preservação do ambiente, à coleta do lixo produzido nas atividades e à reciclagem de inservíveis ; - Possibilitar a expressão dos \textbf{desejos}, saberes e \textbf{afetos}, a \textbf{escuta}, a conversa e a negociação dos conflitos, bem como a amizade entre crianças/crianças, crianças/adultos ; - Promover a autonomia, a autoestima e o desenvolvimento da identidade pessoal, de modo que a \textbf{criança} se sinta pertencente e valorizada quanto ao seu grupo étnico-racial, a sua crença religiosa, a sua cultura, a sua regionalidade, os seus costumes, etc. ; - Fortalecer os vínculos afetivos das \textbf{crianças} tanto nas instituições de Educação \textbf{Infantil} quanto com as suas famílias, de modo a ajudá -las a compreender o mundo e a si mesmas ; - Levar as \textbf{crianças} à refletirem sobre os diversos tipos de preconceito, para promover atitudes de respeito, solidariedade e não-discriminação ; - Incentivar a criação de \textbf{hábitos} de \textbf{higiene} e de preservação do meio ambiente, por meio da coleta do lixo produzido durante as atividades, da reutilização de materiais, da conservação dos recursos naturais, etc. ² As orientações sobre algumas das possibilidades de experiências com \textbf{bebês}, \textbf{crianças} bem \textbf{pequenas} e \textbf{crianças} \textbf{pequenas} irão se repetir nos cinco campos de experiências e estão apresentando uma ideia geral do que poderá ser desenvolvido na educação \textbf{infantil}, no entanto, recomenda -se que seja realizada a adequação da possibilidade, no planejamento do professor, em relação a faixa etária e as condições objetivas da turma atendida. 2.1.5 Observação e avaliação do planejamento A observação e a avaliação do planejamento pedagógico são tão importantes quanto a proposição de experiências significativas.
Assim, os professores de todas as faixas etárias da Educação \textbf{Infantil} precisam assegurar, observar e avaliar no planejamento alguns aspectos, como os seguintes : - Situações de interações e \textbf{brincadeiras} a partir dos objetivos de aprendizagem e desenvolvimento ; - \textbf{Rotinas} que contemplem critérios de regularidade, de continuidade e de diversidade ; - Experiências, espaços, tempo e ambientes que proporcione às \textbf{crianças} o bem-estar, a autonomia, o autocuidado, a descoberta \textbf{sensorial}, a vivência da cultura Mato-Grossense, a tomada de decisões, etc. ; - Diálogo entre crianças/crianças, crianças/adultos para a resolução de conflitos e as manifestações de \textbf{desejos}, sentimentos e opiniões ; - Momentos de \textbf{escuta}, diálogos e acolhimento das famílias ; - Diversidade e regularidade nas estratégias, nos recursos e nos materiais a serem oferecidos às \textbf{crianças}, garantindo -lhes a realização em seu processo de desenvolvimento criador.
Algumas orientações No trabalho com este campo de experiência destaca -se a importância em compor ações pedagógicas capazes de ampliar nas \textbf{crianças} o protagonismo, a autonomia, o conhecimento de si, do outro e do mundo.
Para tanto, a \textbf{criança} necessita conviver com o outro e estabelecer relações, que permitam construir significados, ideias e opiniões, descobrir particularidades sobre si mesma, inclusive que é pertencente a uma família, uma comunidade e uma cultura, despertar o autocuidado, o \textbf{cuidado} com o próximo e a interdependência com o meio, bem como ampliar esse universo pessoal conhecendo outras culturas, identidades e costumes, para respeitar e valorizar a diversidade humana.

% matched lemmas: afeto, apoiar, bebê, brincadeira, conversar, criança, cuidado, desejo, escuta, higiene, hábito, infantil, limpeza, pequeno, rotina, sensorial
\end{textsample}
